% =====================================================================
% Salidas de la consola de Cisco Packet Tracer, tal cual.
%
% Uso: en el preámbulo del documento principal
%     % =====================================================================
% Salidas de la consola de Cisco Packet Tracer, tal cual.
%
% Uso: en el preámbulo del documento principal
%     % =====================================================================
% Salidas de la consola de Cisco Packet Tracer, tal cual.
%
% Uso: en el preámbulo del documento principal
%     % =====================================================================
% Salidas de la consola de Cisco Packet Tracer, tal cual.
%
% Uso: en el preámbulo del documento principal
%     \input{salidas_cisco}
% y en el texto
%     \salidacisco{salidas/orgB/c1_rb4_traceroute_hb2.txt}
%                 {Traceroute de RB4 a HB2}{fig:c1-traceroute}
%     ... como se ve en la Figura~\ref{fig:c1-traceroute} ...
%
% Cada salida está en su propio .txt, copiada literalmente de Packet Tracer:
% LaTeX la lee del fichero, así que se conservan los espacios, la alineación de
% columnas y las líneas en blanco. No hay que copiarla a mano en el .tex.
%
% Fuente: Courier (pcr), monoespaciada como la consola de Packet Tracer.
% Las líneas muy largas (p.ej. los mensajes %OSPF-5-ADJCHG) se parten al
% llegar al borde, igual que hace la consola, sin añadir símbolos.
% =====================================================================
\usepackage[T1]{fontenc}
\usepackage{listings}
\usepackage{xcolor}

% La consola de Packet Tracer tiene 80 columnas. Con \footnotesize, en un
% documento de 11pt con márgenes de 2,5 cm caben ~82 caracteres por línea, así que
% las tablas (show ip ospf neighbor, show ip route) salen sin partirse.
% Si la plantilla final tiene márgenes más grandes, usar [\scriptsize].
\lstdefinestyle{cisco}{
  basicstyle=\fontfamily{pcr}\selectfont\footnotesize,
  columns=fixed,          % todas las letras igual de anchas -> columnas alineadas
  keepspaces=true,        % respeta los espacios tal cual
  showstringspaces=false,
  breaklines=true,        % parte las líneas largas como la consola
  breakatwhitespace=false,
  breakindent=0pt,
  prebreak={}, postbreak={},
  tabsize=8,
  upquote=true,           % comillas rectas, como en la consola
  frame=single,
  rulecolor=\color{black!35},
  framesep=1.5mm,
  xleftmargin=1.5mm, xrightmargin=1.5mm,
  backgroundcolor=\color{white},
  aboveskip=0pt, belowskip=0pt,
}

% \salidacisco[tamaño opcional]{fichero}{pie de figura}{etiqueta}
% Tamaño por defecto \footnotesize (80 columnas caben enteras).
\newcommand{\salidacisco}[4][\footnotesize]{%
  \begin{figure}[htbp]
    \centering
    \lstinputlisting[style=cisco,basicstyle=\fontfamily{pcr}\selectfont#1]{#2}
    \caption{#3}
    \label{#4}
  \end{figure}%
}

% y en el texto
%     \salidacisco{salidas/orgB/c1_rb4_traceroute_hb2.txt}
%                 {Traceroute de RB4 a HB2}{fig:c1-traceroute}
%     ... como se ve en la Figura~\ref{fig:c1-traceroute} ...
%
% Cada salida está en su propio .txt, copiada literalmente de Packet Tracer:
% LaTeX la lee del fichero, así que se conservan los espacios, la alineación de
% columnas y las líneas en blanco. No hay que copiarla a mano en el .tex.
%
% Fuente: Courier (pcr), monoespaciada como la consola de Packet Tracer.
% Las líneas muy largas (p.ej. los mensajes %OSPF-5-ADJCHG) se parten al
% llegar al borde, igual que hace la consola, sin añadir símbolos.
% =====================================================================
\usepackage[T1]{fontenc}
\usepackage{listings}
\usepackage{xcolor}

% La consola de Packet Tracer tiene 80 columnas. Con \footnotesize, en un
% documento de 11pt con márgenes de 2,5 cm caben ~82 caracteres por línea, así que
% las tablas (show ip ospf neighbor, show ip route) salen sin partirse.
% Si la plantilla final tiene márgenes más grandes, usar [\scriptsize].
\lstdefinestyle{cisco}{
  basicstyle=\fontfamily{pcr}\selectfont\footnotesize,
  columns=fixed,          % todas las letras igual de anchas -> columnas alineadas
  keepspaces=true,        % respeta los espacios tal cual
  showstringspaces=false,
  breaklines=true,        % parte las líneas largas como la consola
  breakatwhitespace=false,
  breakindent=0pt,
  prebreak={}, postbreak={},
  tabsize=8,
  upquote=true,           % comillas rectas, como en la consola
  frame=single,
  rulecolor=\color{black!35},
  framesep=1.5mm,
  xleftmargin=1.5mm, xrightmargin=1.5mm,
  backgroundcolor=\color{white},
  aboveskip=0pt, belowskip=0pt,
}

% \salidacisco[tamaño opcional]{fichero}{pie de figura}{etiqueta}
% Tamaño por defecto \footnotesize (80 columnas caben enteras).
\newcommand{\salidacisco}[4][\footnotesize]{%
  \begin{figure}[htbp]
    \centering
    \lstinputlisting[style=cisco,basicstyle=\fontfamily{pcr}\selectfont#1]{#2}
    \caption{#3}
    \label{#4}
  \end{figure}%
}

% y en el texto
%     \salidacisco{salidas/orgB/c1_rb4_traceroute_hb2.txt}
%                 {Traceroute de RB4 a HB2}{fig:c1-traceroute}
%     ... como se ve en la Figura~\ref{fig:c1-traceroute} ...
%
% Cada salida está en su propio .txt, copiada literalmente de Packet Tracer:
% LaTeX la lee del fichero, así que se conservan los espacios, la alineación de
% columnas y las líneas en blanco. No hay que copiarla a mano en el .tex.
%
% Fuente: Courier (pcr), monoespaciada como la consola de Packet Tracer.
% Las líneas muy largas (p.ej. los mensajes %OSPF-5-ADJCHG) se parten al
% llegar al borde, igual que hace la consola, sin añadir símbolos.
% =====================================================================
\usepackage[T1]{fontenc}
\usepackage{listings}
\usepackage{xcolor}

% La consola de Packet Tracer tiene 80 columnas. Con \footnotesize, en un
% documento de 11pt con márgenes de 2,5 cm caben ~82 caracteres por línea, así que
% las tablas (show ip ospf neighbor, show ip route) salen sin partirse.
% Si la plantilla final tiene márgenes más grandes, usar [\scriptsize].
\lstdefinestyle{cisco}{
  basicstyle=\fontfamily{pcr}\selectfont\footnotesize,
  columns=fixed,          % todas las letras igual de anchas -> columnas alineadas
  keepspaces=true,        % respeta los espacios tal cual
  showstringspaces=false,
  breaklines=true,        % parte las líneas largas como la consola
  breakatwhitespace=false,
  breakindent=0pt,
  prebreak={}, postbreak={},
  tabsize=8,
  upquote=true,           % comillas rectas, como en la consola
  frame=single,
  rulecolor=\color{black!35},
  framesep=1.5mm,
  xleftmargin=1.5mm, xrightmargin=1.5mm,
  backgroundcolor=\color{white},
  aboveskip=0pt, belowskip=0pt,
}

% \salidacisco[tamaño opcional]{fichero}{pie de figura}{etiqueta}
% Tamaño por defecto \footnotesize (80 columnas caben enteras).
\newcommand{\salidacisco}[4][\footnotesize]{%
  \begin{figure}[htbp]
    \centering
    \lstinputlisting[style=cisco,basicstyle=\fontfamily{pcr}\selectfont#1]{#2}
    \caption{#3}
    \label{#4}
  \end{figure}%
}

% y en el texto
%     \salidacisco{salidas/orgB/c1_rb4_traceroute_hb2.txt}
%                 {Traceroute de RB4 a HB2}{fig:c1-traceroute}
%     ... como se ve en la Figura~\ref{fig:c1-traceroute} ...
%
% Cada salida está en su propio .txt, copiada literalmente de Packet Tracer:
% LaTeX la lee del fichero, así que se conservan los espacios, la alineación de
% columnas y las líneas en blanco. No hay que copiarla a mano en el .tex.
%
% Fuente: Courier (pcr), monoespaciada como la consola de Packet Tracer.
% Las líneas muy largas (p.ej. los mensajes %OSPF-5-ADJCHG) se parten al
% llegar al borde, igual que hace la consola, sin añadir símbolos.
% =====================================================================
\usepackage[T1]{fontenc}
\usepackage{listings}
\usepackage{xcolor}

% La consola de Packet Tracer tiene 80 columnas. Con \footnotesize, en un
% documento de 11pt con márgenes de 2,5 cm caben ~82 caracteres por línea, así que
% las tablas (show ip ospf neighbor, show ip route) salen sin partirse.
% Si la plantilla final tiene márgenes más grandes, usar [\scriptsize].
\lstdefinestyle{cisco}{
  basicstyle=\fontfamily{pcr}\selectfont\footnotesize,
  columns=fixed,          % todas las letras igual de anchas -> columnas alineadas
  keepspaces=true,        % respeta los espacios tal cual
  showstringspaces=false,
  breaklines=true,        % parte las líneas largas como la consola
  breakatwhitespace=false,
  breakindent=0pt,
  prebreak={}, postbreak={},
  tabsize=8,
  upquote=true,           % comillas rectas, como en la consola
  frame=single,
  rulecolor=\color{black!35},
  framesep=1.5mm,
  xleftmargin=1.5mm, xrightmargin=1.5mm,
  backgroundcolor=\color{white},
  aboveskip=0pt, belowskip=0pt,
}

% \salidacisco[tamaño opcional]{fichero}{pie de figura}{etiqueta}
% Tamaño por defecto \footnotesize (80 columnas caben enteras).
\newcommand{\salidacisco}[4][\footnotesize]{%
  \begin{figure}[htbp]
    \centering
    \lstinputlisting[style=cisco,basicstyle=\fontfamily{pcr}\selectfont#1]{#2}
    \caption{#3}
    \label{#4}
  \end{figure}%
}
